The features presented in this part, collectively, yield a logical framework capable of representing the logical foundations of most formal systems encountered in practice. Our modular approach allows for picking and choosing the required features for a given formalization in a compositional manner, while avoiding the potential overhead of unnecessary rules for a given formalization.\medskip

%We primarily focused on rule-based typing features, such as the components of Martin-Löf type theory. 

\paragraph{Additional Features} Notably, many formal systems offer additional features as ``syntactic sugar'', with an abbreviation-like semantics that can already be fully represented in the underlying logic, but which allow for formalizing content more conveniently. Examples include record and function updates in PVS (e.g. ``$r$ is defined as the record $s$, except for field $a$, which has value $t$ instead''), \emph{let}-operators that introduce new variables as abbreviations for complex expressions, or \emph{sections} in Coq (where local variables are introduced as constants, references to which are lambda-abstracted away when closing a section).

In the interest of preserving the syntactic representation of an imported library as much as possible, we prefer adding these features in the formalization of a system's foundation. Due to their trivial semantics, these features rarely pose a challenge for implementation, and because they can be elaborated away (and usually are in the system itself during checking), they do not introduce hurdles for library integration.

We have implemented some of these convenience features, including a \emph{let}-operator\footnote{\url{https://gl.mathhub.info/MMT/LFX/tree/master/scala/info/kwarc/mmt/LFX/datatypes}}, an object-level substitution function\footnote{ibid.}, a primitive polymorphic type of \emph{lists}\footnote{ibid.}, and Coq-like sections\footnote{\url{https://github.com/UniFormal/MMT/blob/master/src/mmt-coq/src/info/kwarc/mmt/coq/Features.scala}}, the latter of which is very naturally implemented as a structural feature and used in our import of the Coq library~\cite{coqmmt}. Furthermore, Fulya Horozal previously implemented an extension of \LF with \emph{finite sequences} and \emph{flexary} operators~\cite{Horozal:afddl14}.
\medskip

\paragraph{Module-Level Extensions} The \cn{Mod}-operator described in \autoref{sec:lf:records} additionally uses Module-level expressions (namely references to theories), but is still a rule-based object-level typing feature of the underlying logic. However, it is also possible to extend \mmt's module system itself in various ways. These are (regarding their implementation) rather recent additions to the \mmt system and are not central to this thesis, however, they offer interesting extensions for the goal of this part. In particular, module-level features are entirely orthogonal to the logic-level typing feature discussed in this part; hence they are intrinsically compatible with the modular development approach central to this thesis, and composable with all features herein:

\begin{itemize}
	\item Structural features have by now been extended to the module level. Instead of elaborating into a list of declarations, they can instead elaborate into a list of modules.\footnote{As-of-yet unpublished}
	\item \emph{Implicit morphisms} can be used to treat arbitrary theory morphisms in a manner similar to plain includes; in particular they preserve the (fully qualified) names of the symbols in the source theory, whose semantics in the target theory is provided by the morphism. This approach has been published in \cite{RabMue:WADT18}, which also contains a novel description of \mmt's module system and its semantics. While the details therein are largely irrelevant and unnecessarily elaborate for this thesis, it is highly recommended for readers interested in designing module systems.
	\item In this thesis, we treat all theories as named modules with explicitly declared contents. However, \mmt treats module references as \mmt terms, and \emph{theory combinators} (such as unions or pushouts) allow for building complex terms representing anonymous theories. In particular, the theory combinators of \emph{mathscheme}~\cite{mathscheme} have been implemented in \mmt in \ednote{TODO Add citation}.
\end{itemize}

\paragraph{} In addition to allowing the formalization of the foundations of formal systems, our modular logical framework is used as a basis for the Math-in-the-Middle library (see \autoref{sec:lf:mitm}), which serves as a case study and explorative playground for evaluation and inspiration for additional features. 

Additionally, the latter allows us to experiment with combining and unifying various different approaches and combinations of features, e.g. \cn{Mod}-types for theories defined via theory combinators that are the codomains of implicit morphisms, or identifying natural numbers as a \wtypesym-type with a declarative implementation using the Peano axioms, and a separate type populated by literals.